\begin{proofbox}
\textbf{\textcolor{CProof}{方法一：展开平方与立方}}
设 $t\ne0$ 为实数，$u=t+t^{-1}$。
\begin{center}
\begin{tikzpicture}[x=1.9cm,y=1.9cm,font=\small]
\node[font=\bfseries] at (2.7,3.35) {$t=2$：边长 $t+1/t$ 的正方形};
\fill[CStatement!13] (0,0) rectangle (2,2);
\fill[CExample!18] (2,0) rectangle (2.5,2);
\fill[CExample!18] (0,2) rectangle (2,2.5);
\fill[CStatement!28] (2,2) rectangle (2.5,2.5);
\draw (0,0) rectangle (2.5,2.5);
\draw (2,0)--(2,2.5);\draw (0,2)--(2.5,2);
\node at (1,1) {$t^2=4$};\node[rotate=90] at (2.25,1) {$t\cdot t^{-1}=1$};
\node at (1,2.25) {$t\cdot t^{-1}=1$};
\node[anchor=west] at (3,2.3) {$t^{-2}=1/4$};
\draw (2.35,2.3)--(2.8,2.3);
\node at (1,-0.22) {$t=2$};\node at (2.25,-0.22) {$1/t$};
\node[align=left,anchor=west] at (3,1.1) {两个长方形的面积\\各为 $1$；\\除去它们共减 $2$。};
\end{tikzpicture}
\end{center}

图中 $t=2>0$，面积分割显示
\[
u^2=t^2+2+t^{-2}\quad\Longrightarrow\quad t^2+t^{-2}=u^2-2.
\]
对任意非零实数 $t$，即使 $t<0$ 不能直接作边长，同一等式仍由完全平方公式成立。再展开立方：
\[
u^3=t^3+3t+3t^{-1}+t^{-3},
\]
故 $t^3+t^{-3}=u^3-3u$。将二次倒数和再平方，得
\[
t^4+t^{-4}=(u^2-2)^2-2=u^4-4u^2+2.
\]
\textbf{范围的必要性与充分性：}
$u^2-4=(t-t^{-1})^2\ge0$，故 $|u|\ge2$。反过来，任意 $|u|\ge2$ 使 $t^2-ut+1=0$ 的判别式 $u^2-4\ge0$；两根乘积为 $1$，均非零，故每个这样的 $u$ 都能取得。边界 $u=\pm2$ 分别对应 $t=\pm1$；$t>0$ 时 $u\ge2$，$t<0$ 时 $u\le-2$。
\end{proofbox}
\clearpage
\begin{proofbox}
\textbf{\textcolor{CProof}{方法二：乘后减去多余的一组}}
\begin{center}
\begin{tikzpicture}[font=\small,x=1cm,y=1cm]
\node[font=\bfseries] at (5.4,3.7) {一次倒数和乘二次倒数和：交叉格又组成 $u$};
\node at (4,2.8) {$t^2$};\node at (8,2.8) {$t^{-2}$};
\node at (1.2,1.8) {$t$};\node at (1.2,0.6) {$t^{-1}$};
\fill[CExample!16] (6,1.2) rectangle (10,2.4);
\fill[CExample!16] (2,0) rectangle (6,1.2);
\draw (2,0) rectangle (10,2.4);\draw (6,0)--(6,2.4);\draw (2,1.2)--(10,1.2);
\node at (4,1.8) {$t^3$};\node at (8,1.8) {$t^{-1}$};
\node at (4,.6) {$t$};\node at (8,.6) {$t^{-3}$};
\node at (5.4,-.7) {绿色两格的和为 $u$；剩下两格的和为 $S_3$。};
\end{tikzpicture}
\end{center}

乘法格直接给出
\[
u(t^2+t^{-2})=(t^3+t^{-3})+(t+t^{-1}).
\]
代入方法一的 $t^2+t^{-2}=u^2-2$，减去绿色交叉格的和 $u$，得
\[
\boxed{t^3+t^{-3}=u(u^2-2)-u=u^3-3u.}
\]
四次同样不用展开四次二项式，只需对二次整体平方，再扣去两个交叉项：
\[
\boxed{t^4+t^{-4}=(t^2+t^{-2})^2-2=(u^2-2)^2-2.}
\]
本法利用已知的低次整体，适合快速推式；$t\ne0$ 与 $|u|\ge2$ 的条件仍须保留。
\end{proofbox}
\clearpage
\begin{proofbox}
\textbf{\textcolor{CProof}{方法三：首尾对称四次方程降次}}
设 $a,b\in\mathbb R$。方程 $t^4+at^3+bt^2+at+1=0$ 的常数项为 $1$，故 $t=0$ 不是根，可以同除 $t^2$。
\begin{center}
\begin{tikzpicture}[font=\small,x=1cm,y=1cm]
\node[font=\bfseries] at (5.5,3.7) {除以 $t^2$ 后，幂次关于 $0$ 成对};
\draw[-Stealth] (.5,0)--(11,0) node[anchor=west] {幂次};
\foreach \x/\k/\v in {1/-2/{$t^{-2}$},3.25/-1/{$at^{-1}$},5.5/0/{$b$},7.75/1/{$at$},10/2/{$t^2$}} {
\draw (\x,-.1)--(\x,.1);\node[anchor=north] at (\x,-.2) {$\k$};\node at (\x,.55) {\v};}
\draw[CStatement,thick] (1,1)--(1,2.7)--(10,2.7)--(10,1);
\draw[CExample,thick] (3.25,1)--(3.25,1.65)--(7.75,1.65)--(7.75,1);
\node[fill=BGProof,text=CStatement] at (5.5,2.7) {$t^2+t^{-2}=u^2-2$};
\node[fill=BGProof,text=CExample] at (5.5,1.65) {$a(t+t^{-1})=au$};
\end{tikzpicture}
\end{center}

图中两对正负幂次各自合并，中央保留常数 $b$，于是
\[
(t^2+t^{-2})+a(t+t^{-1})+b=0
\]
化为
\[
\boxed{u^2+au+b-2=0,\qquad |u|\ge2.}
\]
先求这个二次方程的根，舍去 $(-2,2)$ 内的值；对剩下的每个 $u$ 解
\[
t^2-ut+1=0,\qquad t=\frac{u\pm\sqrt{u^2-4}}2.
\]
当 $|u|>2$ 时对应两个互为倒数的不同实根；当 $u=2$ 或 $-2$ 时，两个回代值重合为 $t=1$ 或 $-1$。全部变形在 $t\ne0$ 下等价，因此回代后不会漏解或引入新解。
\end{proofbox}
\clearpage
\begin{proofbox}
\textbf{\textcolor{CProof}{方法四：指数移动得到通用递推}}
记 $S_n=t^n+t^{-n}$（整数 $n\ge0$），$S_0=2$、$S_1=u$。
\begin{center}
\begin{tikzpicture}[font=\small,x=1cm,y=1cm]
\node[font=\bfseries] at (5.5,4.2) {乘 $t$ 指数加 $1$；乘 $t^{-1}$ 指数减 $1$};
\draw[-Stealth] (.3,0)--(10.7,0) node[anchor=west] {指数};
\foreach \x/\k in {1/{-n-1},2.5/{-n},4/{-n+1},7/{n-1},8.5/{n},10/{n+1}} {
\draw (\x,-.1)--(\x,.1);\node[anchor=north] at (\x,-.25) {$\k$};}
\foreach \x in {2.5,8.5} {\fill[CStatement] (\x,0) circle (2pt);}
\draw[-Stealth,CExample,thick] (2.5,.25) to[bend left=35] (4,.25);
\draw[-Stealth,CExample,thick] (8.5,.25) to[bend left=35] (10,.25);
\draw[-Stealth,CTrap,thick,dashed] (2.5,.25) to[bend right=35] (1,.25);
\draw[-Stealth,CTrap,thick,dashed] (8.5,.25) to[bend right=35] (7,.25);
\node at (5.5,1.45) {绿色实线：乘 $t$；红色虚线：乘 $t^{-1}$};
\draw[CStatement,thick] (1,1.95)--(1,2.9)--(10,2.9)--(10,1.95);
\node[fill=BGProof] at (5.5,2.9) {两端组成 $S_{n+1}$};
\node at (5.5,3.55) {内侧 $n-1$ 与 $-n+1$ 组成 $S_{n-1}$};
\node at (5.5,-1.05) {位置为指数示意，中央省略间隔；$n=1$ 时内侧两项都为 $1$。};
\end{tikzpicture}
\end{center}

按图中指数的移动，乘积的四项恰好组成两组：
\[
\begin{aligned}
uS_n&=(t+t^{-1})(t^n+t^{-n})\\
&=t^{n+1}+t^{n-1}+t^{-n+1}+t^{-n-1}\\
&=S_{n+1}+S_{n-1}\qquad(n\ge1).
\end{aligned}
\]
所以 $\boxed{S_{n+1}=uS_n-S_{n-1}}$。依次算得
\[
S_2=u^2-2,\quad S_3=u^3-3u,\quad S_4=u^4-4u^2+2.
\]
从常数 $S_0$ 与一次式 $S_1$ 出发，每次仅乘 $u$ 并减去已知多项式，故所有 $S_n$ 都是 $u$ 的多项式。$n=1$ 时两项 $t^0$ 均须保留，因此减的是 $S_0=2$。
\end{proofbox}
\clearpage
\begin{proofbox}
\textbf{\textcolor{CProof}{补充：倒数差的恒等式与范围}}
令 $v=t-t^{-1}$，则展开得
\[
v^2=t^2-2+t^{-2},\qquad v^3=t^3-3t+3t^{-1}-t^{-3},
\]
即 $t^2+t^{-2}=v^2+2$、$t^3-t^{-3}=v^3+3v$。
\begin{center}
\begin{tikzpicture}[font=\small]
\begin{axis}[at={(0,0)},width=6.3cm,height=6cm,axis lines=middle,
 xmin=-3.5,xmax=3.5,ymin=-4.5,ymax=4.5,xtick={-1,1},ytick={-2,2},
 yticklabel style={anchor=west,xshift=4pt},
 xlabel={$t$},title={和型 $u=t+t^{-1}$}]
\addplot[CStatement,very thick,domain=.24:3.4,samples=130] {x+1/x};
\addplot[CStatement,very thick,domain=-3.4:-.24,samples=130] {x+1/x};
\addplot[only marks,CStatement] coordinates {(1,2) (-1,-2)};
\end{axis}
\begin{axis}[at={(6.5cm,0)},width=6.3cm,height=6cm,axis lines=middle,
 xmin=-3.5,xmax=3.5,ymin=-4.5,ymax=4.5,xtick=\empty,ytick=\empty,
 xlabel={$t$},title={差型 $v=t-t^{-1}$}]
\addplot[CExample,very thick,domain=.22:3.4,samples=130] {x-1/x};
\addplot[CExample,very thick,domain=-3.4:-.22,samples=130] {x-1/x};
\addplot[only marks,CExample] coordinates {(1,0) (-1,0)};
\end{axis}
\end{tikzpicture}
\end{center}

两幅图的坐标范围与尺度相同。左图正支最低为 $2$、负支最高为 $-2$；右图每条分支从一端趋向另一端，均覆盖全部纵坐标。

严格地说，任给实数 $v$，回代方程 $t^2-vt-1=0$ 的判别式为 $v^2+4>0$，两根乘积为 $-1$，均非零，故 $v$ 的范围确为 $\mathbb R$。
\end{proofbox}
\clearpage

\begin{proofbox}
\textbf{\textcolor{CProof}{补充推导：三角函数的和、积如何整体换元}}
\textbf{适用条件：} $0\le\theta\le\pi/2$，设
$t=\sin\theta+\cos\theta$。这里的 $t$ 与前面倒数型题目的原变量无关。

\textbf{第一步（平方，不跳项）：}
\[
\begin{aligned}
t^2&=(\sin\theta+\cos\theta)^2\\
&=\sin^2\theta+2\sin\theta\cos\theta+\cos^2\theta\\
&=1+2\sin\theta\cos\theta.
\end{aligned}
\]
把式子移项并除以 $2$：
\[
\boxed{\sin\theta\cos\theta=\frac{t^2-1}{2}}.
\]
\textbf{第二步（先展开目标式，再代换）：}
\[
\begin{aligned}
(1+\sin\theta)(1+\cos\theta)
&=1+\sin\theta+\cos\theta+\sin\theta\cos\theta\\
&=1+t+\frac{t^2-1}{2}\\
&=\frac{t^2+2t+1}{2}
=\boxed{\frac{(t+1)^2}{2}}.
\end{aligned}
\]
\textbf{第三步（证明范围和等号）：}
由于 $\sin\theta,\cos\theta\ge0$，有
$t^2=1+2\sin\theta\cos\theta\ge1$，且 $t\ge0$，故 $t\ge1$。
另一方面，
\[
0\le(\sin\theta-\cos\theta)^2
=1-2\sin\theta\cos\theta
\]
说明 $2\sin\theta\cos\theta\le1$，于是 $t^2\le2$，得到
\[
\boxed{1\le t\le\sqrt2}.
\]
$t=\sqrt2$ 当且仅当 $\sin\theta=\cos\theta$；在给定角度范围内即
$\theta=\pi/4$。在这个区间 $t+1>0$，故 $(t+1)^2/2$
随 $t$ 增大而增大，最大值为 $\frac{(1+\sqrt2)^2}{2}$。

\textbf{关键区别：} 倒数型利用两项\emph{乘积固定为 $1$}；本题利用两项
\emph{平方和固定为 $1$}。共同思路是先识别重复出现的整体，再检查新变量范围。
\end{proofbox}
\clearpage
